\begin{figure}[htbp]
\centering
\begin{tikzpicture}
\begin{axis}[width=.87\textwidth,height=60mm,axis lines=left,
 xmin=0,xmax=2.15,ymin=0,ymax=1.15,xlabel={$x$},ylabel={$f(x)$},
 samples=120]
\pgfplotsinvokeforeach{0,...,7}{
 \pgfmathsetmacro{\leftx}{#1/4}
 \pgfmathsetmacro{\rightx}{(#1+1)/4}
 \pgfmathsetmacro{\rectheight}{exp(-#1/4)}
 \edef\rectanglecmd{\noexpand\draw[draw=arkivblaa!65,fill=arkivblaa!13]
   (axis cs:\leftx,0) rectangle (axis cs:\rightx,\rectheight);}
 \rectanglecmd
}
\addplot[arkivrod,thick,domain=0:2]{exp(-x)};
\node[anchor=north east,font=\small] at (rel axis cs:.96,.96)
 {$f(x)=\e^{-x}$, $N=8$};
\end{axis}
\end{tikzpicture}
\caption{Venstre Riemannsum for en monoton aftagende, glat funktion på
$[0,2]$. Hvert rektangel ligger over den tilhørende grafdel, så summen er
en øvre approksimation. Højre prøvepunkter ville give en nedre sum.
Inddelingens bredde skal gå mod nul i definitionen af integralet;
figuren viser én endelig inddeling i dimensionsløse variable.}
\label{fig:integral-riemann}
\end{figure}
